\cotefile{84}
\scanpage{67}
% lines 1–3, opening
Notons que pour un \unsure{cert.|ext.} cas donné, il revient au $=$ de se donner $T$ tel que $\psi(T)=\unsure{\chi|\gamma}$, \unsure{où|ou} un ``hom de $\chi$-relèvement''
% lines 4–7, upper centre and right
\[
\underline{O}\xrightarrow{\psi'\,(=\psi'_T)}E
\]
\margin{(néc. de la forme $\lambda\mathbin{\unsure{-|\mapsto}}\lambda T$, où $\underset{\substack{\shortparallel\\\psi'(1)}}{T}\in\Gamma E$)}
tel que
\[
\psi\psi'(\deleted{\unsure{=}}\lambda)=\chi\lambda
\]
(ce qui équivaut en effet à $\psi\psi'(1)=\chi$ i.e. $\psi(T)=\chi$).
% lines 8–11, above centre
Il revient aussi au $=$ de se donner un hom
\[
E\xrightarrow{\pi\,(=\pi_T)}\unsure{\underline M|\underline L|M}
\]
qui soit une $\chi$-ré\deleted{\unsure{t}}traction p.r. à $i$, i.e. tel que
\[
\pi\circ i=\chi\,\mathrm{id}_{\underline M}
\]
% lines 12–14, centre
La relation entre $\pi_{\deleted{\unsure{T}}}$ et $T$ est la suivante
\[
\begin{aligned}
i(\pi(x))&=\underset{\substack{\shortparallel\\\psi(T)}}{\chi}\,x-T\,\psi(x)
\qquad\text{i.e.}\quad\pi=\chi\,\mathrm{id}-\psi\otimes T\\
(&=\psi(T)x-T\psi(x)\overset{\mathrm{def}}{=}x\mathbin{\underset{\psi}{\wedge}}T)
\end{aligned}
\]
% lines 15–21, below centre
On aura, pour $T$ \unsure{fixée|fixé} telle que $\psi(T)=\chi$, et pour tout $x$
\[
\begin{aligned}
\deleted{\unsure{\psi}}\psi(\pi(x))&=\chi\,\psi(x)-\psi(T)\,\psi(x)=0\qquad\text{car }\psi(T)=\chi\\
&\text{i.e.}\quad\pi(x)\in i(L)
\end{aligned}
\]
et si $x\in i(L)$ i.e. $\psi(x)=0$, on a bien
\[
\unsure{i}\bigl(\pi(x)\bigr)=\chi\,x\qquad \deleted{\text{i.e. }i\pi(x)=\chi}
\]
% lines 22–27, foot
Inversement, si $\pi$ est donnée, considérons
\[
x\longmapsto\chi\deleted{\unsure{a}}x-\pi x\quad:\quad E\longrightarrow\deleted{M}E
\]
elle est nulle sur $\underline M$, donc se factorise en
\[
\begin{aligned}
E/\unsure{\underline M|L|M}\simeq O&\xrightarrow{\psi'}E\\
\lambda&\longmapsto\lambda T
\end{aligned}
\]
cette flèche $\psi'$ donnée par un $T\in E$, et on aura
\[
\chi\deleted{\unsure{a}}x-\pi x=\psi(x)T\qquad\text{i.e.}\quad\boxed{\pi(x)=\chi\deleted{\unsure{a}}x-\psi(x)T}
\]
% archive sheet number: 66
\scanpage{68}
% lines 1–5, opening
Prenant le $\psi$ des deux \unsure{mbs|mbrs}, et comme $\pi(x)\in L$ donc $\psi(\pi(x))=0$, on trouve
\[
\unsure{\chi|\gamma}\,\psi(x)-\psi(x)\,\psi(T)=0\qquad\text{i.e.}\quad\psi(x)(\chi-\psi(T))=0
\]
pour tte section locale $x$ de $O$. Prenant $x$ telle que $\psi(x)=1$, on trouve $\chi-\psi(T)=0$ i.e. $\psi(T)=\chi$.
% lines 6–8, two given sequences
Soient
\[
\begin{aligned}
0\mathrel{\text{---}} L_1\xrightarrow{\unsure{i'_1|i_1}}E_1\xrightarrow{\psi_1}O\longrightarrow0
&\qquad T_1\in E_1\qquad\psi_1(T_1)=\chi\\
0\mathrel{\text{---}}L_2\xrightarrow{i_2}E_2\overset{\psi_2}{\mathrel{\text{---}}}O\mathrel{\text{---}}0
&\qquad T_2\in E_2\qquad\psi_2(T_{\unsure{1|2}})=\chi
\end{aligned}
\]
% lines 9–11, proposed sequence
on veut définir
\[
0\longrightarrow L_1\otimes L_2\xrightarrow{i}
\underset{\substack{\shortparallel\,\mathrm{def}\\(E_1,T_1)\mathbin{\underset{\chi}{\wedge}}(E_2,T_2)}}{E}
\xrightarrow{\psi}\underline O\longrightarrow0,\qquad T\in E\quad\psi(T)=\chi
\]
% lines 12–15, middle diamond
Considérons
\diagram{\begin{tikzcd}[ampersand replacement=\&,column sep=small]
\& L_1\otimes L_2\arrow[dl]\arrow[dr,dash]\&\\
L_1\otimes E_2\arrow[dr]\&\& E_{1}\otimes L_2\arrow[dl]\\
\&\deleted{\unsure{O}}E_1\otimes E_2\&, 
\end{tikzcd}}
% lines 16–18, filtration and its intervals
qui correspond à une suite de comp. \unsure{sur|de|dans} $E_1\otimes E_2$
\[
0\mathrel{\underset{L_1\otimes L_2}{\underbrace{\ \subset\ }}}L_1\otimes L_{\unsure{1|2}}
\mathrel{\underset{L_1\oplus L_2}{\underbrace{\ \subset\ }}}L_1\otimes E_2+E_1\otimes L_2
\mathrel{\underset{O}{\underbrace{\ \subset\ }}} E_1\otimes E_2,
\]
% lines 19–22, foot and linked note
On a sur $E_1\otimes E_2$ une structure d'ext.
\[
0\longrightarrow
\underset{\substack{\shortparallel\,\mathrm{déf}\\ L_1\otimes E_2\mathbin{\underset{L_1\otimes L_2}{+}}E_1\otimes L_2}}{E_1*_{O}E_2}
\longrightarrow E_1\otimes E_2\longrightarrow O\longrightarrow0
\]
\margin{$\longleftrightarrow$ NB dans la notation $E_1,E_2$, la structure d'ext. est sous-entendue}
% archive sheet number: 67
\scanpage{69}
% lines 1–2, opening
avec
\[
0\to L_1\otimes L_2\longrightarrow E_1*E_2\longrightarrow\deleted{E}L_1\oplus\deleted{E}L_2\longrightarrow0
\]
% lines 3–6, canonical map
On va définir un hom can
\[
E_1*E_2\xrightarrow{\mu=\mu_{E_1,E_2}}L_1\otimes L_2
\]
i.e. un diagramme commutatif d'hom linéaires
% lines 7–9, diamond
\diagram{\begin{tikzcd}[ampersand replacement=\&,column sep=small]
\&L_1\otimes L_2\arrow[dl]\arrow[dr]\&\\
L_1\otimes E_2\arrow[dr,"\alpha_1"']\&\deleted{\unsure{\to|-}}\&E_{1}\otimes L_2\arrow[dl,"\deleted{\unsure{\beta}}\alpha_2"]\\
\&L_1\otimes L_2\&
\end{tikzcd}}
% lines 10–12, two prescriptions
On définit
\[
\alpha\unsure{=|_1=}\mathrm{id}_{L_1}\otimes\pi_{T_2}\qquad\qquad\beta=\pi_{T_1}\otimes\mathrm{id}_{L_2}
\]
i.e.
% lines 13–16, evaluation of the prescriptions
\[
\begin{aligned}
\alpha_1(\underset{\substack{\rotatebox{90}{$\in$}\\\Gamma L_1}}{\deleted{\unsure{x}}u_1}\otimes
\underset{\substack{\rotatebox{90}{$\in$}\\\Gamma E_2}}{x_2})
&=u_1\otimes\pi_{T_2}(x_2)=u_1\otimes(\unsure{\chi|\gamma}x_2-T_2\psi(x_2))\\
\deleted{\unsure{\beta_2|\beta}}\alpha_2(\underset{\substack{\rotatebox{90}{$\in$}\\\Gamma E_1}}{x_1}\otimes
\underset{\substack{\rotatebox{90}{$\in$}\\\Gamma L_2}}{u_2})
&=\pi_{T_1}(x_1)\otimes u_2=(\chi x_1-T_1\inserted{\psi(x_1)})\otimes u_2
\end{aligned}
\]
% lines 17–18, common restriction and right-hand note
on a
\[
\begin{aligned}
\alpha_1(\deleted{\unsure{x}}u_1\otimes u_2)&=\chi\,u_1\otimes u_2\\
\alpha_2(u_1\otimes u_2)&=\chi\,u_1\otimes u_2
\end{aligned}
\]
\margin{OK !}
% lines 19–23, foot
On posera
\[
\begin{aligned}
E_1\wedge E_2={}&\text{extension de }\underline O_{\unsure{X|x}}\text{ par }L_1\otimes L_2\\
&\text{déduite de l'ext. de }O\text{ par}\\
&E_1\mathbin{\unsure{\otimes|*}}E_2\text{ grâce à}\\
&\mu:E_1*E_2\mathbin{\unsure{-|\to}}\unsure{E_1|L_1}\otimes L_2\quad\text{ci-dessus}
\end{aligned}
\]
% archive sheet number: 68
\scanpage{70}
% lines 1–4, upper rectangle
On a donc un hom de suites exactes
\diagram{\begin{tikzcd}[ampersand replacement=\&,column sep=small]
0\arrow[r]\&E_1*_{O}E_2\arrow[r,"i"]\arrow[d,"{\deleted{\unsure{\mu=\mu_{E_1,E_2}}}}"']\&
E_1\otimes E_2\arrow[r,"{x_1\otimes x_2\mathbin{\unsure{\to|-}}\psi_1(x_1)\psi_2(x_2)}"]\arrow[d,"\unsure{\mu}"']\&
O\arrow[r]\arrow[d,"\mathrm{id}"]\&0\\
0\arrow[r]\&\deleted{E}L_1\otimes L_2\arrow[r,"i"]\&E_1*E_2\arrow[r,"\psi"]\&O\arrow[r]\&0
\end{tikzcd}}
% lines 5–8, splitting
Il faut définir un $\unsure{\chi|\gamma}$-scindage de l'ext. $E_1\wedge E_2$.

Il revient au $=$ \deleted{\unsure{$T$ de $E_1\wedge E_2$|$E_1*E_2$}} de se donner
\[
E_1\mathbin{*}E_2\longrightarrow L_1\otimes L_2
\]
qui sur $L_1\otimes L_2$ se réduise à la mult. par $\chi$.
% lines 9–13, pushout and mapping property
Comme
\[
E_1*E_2=E_1\otimes E_2\mathbin{\underset{\deleted{\unsure{E_1\otimes E_2}}E_1*_{O}E_2}{+}}
(\deleted{E}L_1\otimes L_2)
\]
se donner un hom lin. de $E_1*E_2$ dans un module $M$ revient au $=$ qu'un diagramme comm
% lines 14–18, centre diagram
\diagram{\begin{tikzcd}[ampersand replacement=\&,column sep=small]
\&E_1*_{O}E_2\arrow[dl,"\mathrm{inc}"']\arrow[dr,"\mu"]\&\\
E_1\otimes E_2\arrow[dr,dashed]\arrow[ddr,"\rho=\pi"']\&\&L_1\otimes L_2\arrow[dl,dashed]\arrow[ddl,"{\sigma=\chi\,\mathrm{id}_{L_1\otimes L_2}}"]\\
\&E_1\mathbin{\unsure{*|\wedge}}E_2\arrow[d,"\pi_T"]\&\\
\&M(=E_1\otimes E_2)\quad )\&
\end{tikzcd}}
% lines 19–23, last paragraph
\deleted{\unsure{Et}} on a ici $M=L_1\otimes L_2$. Considérons
\[
\rho:E_1\otimes E_2\longrightarrow L_1\otimes L_2\qquad \rho=\pi_{E_1}\otimes\pi_{E_2}
\]
i.e.
\[
\rho(x_1\otimes x_2)=
\underbrace{(\chi x_1-\psi_1(x_1)T_1)}_{\pi_{T_1}(x_1)}\otimes
\underbrace{\bigl(\chi(x_2)-\psi_2(x_2)T_2\bigr)}_{\pi_{T_2}(\unsure{x_2|x_1})}
\]
Si $\overset{\inserted{=u_1}}{x_1}\in L_1$, on trouve \deleted{$\unsure{\chi u_1}$} (comme $\psi_1(x_1)=0$)
% archive sheet number: 69
\scanpage{71}
% lines 1–3, opening computations
\[
\begin{aligned}
&\deleted{\rho(\overbrace{u_1\otimes x_2+x_1\otimes u_2}^{\unsure{\in} E_1*E_2})=\unsure{\chi|\gamma}(u_1\otimes\chi(x_2}\\
&\rho(u_1\otimes x_2)=\chi u_1\otimes\bigl(\chi\deleted{(}x_2\deleted{)}-\psi_2(x_2)T_2\bigr)
\end{aligned}
\]
Mais
\[
\mu_{\deleted{\unsure{E_2|E_1,E_2}}}(u_1\otimes x_2)\overset{\mathrm{def}}{=}u_1\otimes
\underbrace{(\chi\deleted{\unsure{x}}x_2-\psi_2(x_2)T_2)}_{\unsure{\in}}
\]
% lines 4–6, restriction and diagonal note at left
\unsure{On|il faut} prendra donc $\sigma=\chi\,\mathrm{id}_{\deleted{E}L_1\otimes L_2}^{*}$ pour avoir commutativité du côté $L_1\otimes E_2\subset E_1*E_2$, et de $=$ pour $E_1\otimes L_2$. On gagne !
% slanted
\margin{* C'est ce qu'il fallait prendre \unsure{a|à} priori, si on voulait que $\pi_T$ soit une $\chi$-\unsure{trivialisation|trivialis.}}
% lines 7–13, regular scalar
Si $\chi$ est regulier sur \deleted{$\unsure{L_1,L_2}$} $L_1\otimes L_2$ (p.ex. $\chi$ régulier \unsure{sur|dans} $O$ et $L_1,L_2$ plats sur $O$, donc $L_1\otimes L_2$ plat aussi...) alors on peut définir $T$ par
\[
\begin{aligned}
\chi T&=\deleted{\unsure{\rho}}\mu'(T_1\otimes T_2)\overset{\mathrm{déf}}{=}T_1\deleted{\otimes}\unsure{*|\wedge}T_2\\
\text{i.e.}\qquad T&=\frac{1}{\chi}T_1\deleted{\otimes}*T_2\qquad(!)
\end{aligned}
\]
% lines 14–17, projection and its rearrangement
\deleted{Considérons} En effet considérons \deleted{$\unsure{\pi_T}$}
\[
\begin{aligned}
\pi_T:\ E&\longrightarrow L_1\otimes L_2\\
x&\longmapsto\chi x-T\psi(x)\qquad\text{d'où}\qquad T\psi(x)=
\deleted{\unsure{\chi\pi_T(x)}}\inserted{\chi x-\pi_T(x)}
\end{aligned}
\]
% lines 18–21, heavily struck calculation at left and continuation at right
\deleted{\unsure{puis|posons} faisons $x=T$, on trouve, comme $\psi(T)=\chi$}
\[
\begin{aligned}
&\deleted{\chi\,\pi_T=\unsure{\chi} \mathbin{\mapsto}\chi^2x-\chi T\psi(x)}\\
&\deleted{\chi T=\chi T\unsure{.}}
\end{aligned}
\]
Faisons $x=T_1*T_2$, d'où
\[
\begin{aligned}
\psi(x)&=\deleted{\unsure{\chi}}\psi(T_1)\cdot\psi(T_2)=\chi^2\\
\text{d'où}\qquad\chi^2T&=\chi\,T_1*T_2-
\underbrace{\pi_T(T_1*T_2)}_{\pi_{T_1}(\underset{\shortparallel\,0}{T_1})\otimes\pi_{T_{\unsure{1|2}}}(\underset{\shortparallel\,0}{T_2})}
\end{aligned}
\]
% lines 22–25, diagonally struck lower paragraph and display
\deleted{Soient \unsure{$x_1,x_2$} $x_1\in\GammaP_1$, $x_2\in\GammaP_2$, tels q. $\psi_1(x_1)=1$, $\psi_2(x_2)=1$ prenons $x=x_1*x_2$, on aura}
\[
\deleted{T=\chi x_1*x_2-\underbrace{\pi_T(x_1*x_2)}_{(\chi x_1-T_1)(\chi x_2-T_2)}}
\]
% archive sheet number: 70
\scanpage{72}
% lines 1–5, diagonally ruled attempt and right-hand continuation
\deleted{Faisons $\unsure{\to}$}
\[
\deleted{x_1=\frac{1}{\unsure{\chi|\gamma}_1}T_1\qquad x_2=\frac{1}{\chi}T_2\qquad(!)}
\]
\deleted{on trouve}
\[
\deleted{T=\chi\underbrace{\bigl(\tfrac1\chi T_1\bigr)*\bigl(\tfrac1\chi T_2\bigr)}_{\frac1\chi T_1*T_2}-}
\]
\margin{Faisons $x_{1}=T_1*T_2$, d'où}
% lines 6–8, cancellation
d'où
\[
\chi^2T=\chi\,T_1*T_2
\]
d'où en divisant par $\chi$
\[
\boxed{\chi T=T_1*T_2}
\]
% lines 9–13, induced application
\uline{NB} L'application
\[
\begin{aligned}
E_1\times E_2&\mathbin{-}E_1\otimes E_2\longrightarrow E_1*E_2\\
x_1,x_2&\longmapsto x_1\otimes x_2,\ \ x_1\otimes x_2\longmapsto x_1*x_2
\end{aligned}
\]
applique
\[
P_1\times P_2\longrightarrow P=\psi^{-1}(1)
\]
% lines 14–19, elements and product expansion
On aura, si
\[
\begin{aligned}
x_1&\in\Gamma P_1,\ u_1\in\Gamma L_1\\
x_2&\in\Gamma P_2,\ u_2\in\Gamma L_2
\end{aligned}
\]
\[
\begin{aligned}
(\underset{\substack{\rotatebox{90}{$\in$}\\\deleted{\unsure{\Gamma P_1}}}}{x_1}
+\underset{\substack{\rotatebox{90}{$\in$}\\\Gamma L_1}}{u_1})*
(\underset{\substack{\rotatebox{90}{$\in$}\\\Gamma P_2}}{x_2}
+\underset{\substack{\rotatebox{90}{$\in$}\\\Gamma L_2}}{u_2})
&=\underbrace{x_1*x_2}_{\Gamma P}
+\underbrace{u_1*x_2+x_1*u_2}_{\Gamma(\deleted{E}L_1\otimes L_2)}
+\underbrace{u_1*u_2}_{\shortparallel\,\chi u_1\otimes u_2}\\
&=x_1*x_2+\bigl[u_1\otimes(\overbrace{\chi x_2-\deleted{\unsure{\psi_2(x_2)|\psi(x_2)}}T_2}^{-b_2})
+(\overbrace{\chi x_1-\deleted{\unsure{\psi_1(x_1)|\psi(x_1)}}T_1}^{-b_1})\otimes u_2\\
&\hspace{9em}+\chi\,u_1\otimes u_2\bigr]
\end{aligned}
\]
% lines 20–23, lower left coordinates, boxed calculation and right-edge note
\margin{\uline{NB} $x_1=(0,1)\quad T_1=(b_1,\chi)$\\ $\chi x_1=(0,\chi)$\\ \deleted{$\chi x_1-T_1$}\\ $\chi x_1-T_1=(-b_1,0)$}
\[
\boxed{\begin{aligned}
(x_1+u_1)*(x_2+u_2)&=x_1*x_2+(-u_1\otimes b_2-b_1\otimes u_2+\chi\,u_1\otimes u_2)\\
&=x_1*x_2+(-u_1\otimes b_2\mathbin{\unsure{-|+}}b_1\otimes u_2+\chi\,u_1\otimes u_2)
\end{aligned}}
\]
\margin{\unsure{ces signes vont s'éliminer si} \unsure{subst.}}
% lines 24–27, slanted lower-left note and baseline
% slanted
\margin{\uline{NB} Si $\{b_1\}$ est une base de $L_1$, alors $x_2\longrightarrow x_1*x_2$ est un iso $P_2\simeq P_1*P_2$ de torseurs, compatible \unsure{à|avec} l'isom $u_2\longmapsto -b_1\otimes u_2$ de $L_2$ sur $L_1\otimes L_2\simeq O\otimes L_2=L_2$}
$b_1,b_2$ sont les éléments de $\deleted{E}L_1,L_2$ qui définissent $T_1,T_2$ en termes des $x_1,x_2$
\[
T_1=\chi x_1+b_1,\quad T_2=\chi x_2+b_2
\]
% archive sheet number: 71
\scanpage{73}
% lines 1–3, opening identification and formula
identifiant $E$ à $\deleted{\unsure{E}}L_1\otimes L_2\oplus\underline O_X$ grâce au scindage \uline{$x_1*x_2$}, on trouve, par un calcul immédiat$^{(*)}$
\[
T=\deleted{\unsure{\chi(x_1*x_2)}}\unsure{\chi|\gamma}(x_1*x_2)\mathbin{\unsure{-|+}}\deleted{\unsure{\chi}}b_1\otimes b_2
\]
% lines 3–7, two annotations below the identification
\margin{* NB\\ $\pi_T(x_1*x_2)=\chi\,x_1*x_2-T$ \deleted{\unsure{$\chi$}}\\ donc $T=\chi\,x_1\otimes x_2-\pi_T(x_1*x_2)$\\ Or $\pi_T(x_1*x_2)=\pi_{T_1}(\deleted{\unsure{u}}x_1)\otimes\pi_{T_2}(x_2)$\\ $\forall\,x_1\in E_1,\ x_2\in E_2$}
\margin{Re NB : on \unsure{vérifie|vérifierait} aussitôt que ça ne change \unsure{pas} si on remplace $x_1$ par $x_1+u_1$ ou $x_2$ par $x_2+u_2$ ($u_1\in\Gamma L'_1,\ u_2\in\Gamma L_2$)\\ Attention signe !}
% lines 8–12, three braced definitions
Pour éliminer le signe, il faudrait poser
\[
\left\{\begin{aligned}
b_{T_1,x_1}&=-b_1&\quad\text{i.e.}\quad b_{T_1,x_1}&=\chi x_1-T_1=\pi_{T_1}(x_1)\\
b_{T_2,x_2}&=-b_2&\text{i.e.}\quad b_{T_2,x_2}&=\chi x_2-T_2=\pi_{T_2}(x_{\unsure{1|2}})\\
b_{T,x_1*x_2}&=-b&\text{i.e.}\qquad b&=\chi\,x_1\mathbin{\unsure{*|\otimes}}x_2-T=\pi_T(x)\\
&&&(x=x_1\otimes x_2)
\end{aligned}\right.
\]
% lines 13–17, consequence and formula for T
et on trouvera
\[
\deleted{\unsure{b_2|b}}\ b=-b_{T,x_1*x_2}=-b_1\otimes b_2=-b_{T_1,x_1}\otimes b_{T_2,x_2}
\]
d'où
\[
\boxed{b_{T,x_1*x_2}=b_{T_1,x_1}\otimes b_{T_2,x_2}}
\]

On trouve alors
\[
T=\chi\cdot(x_1*x_2)-\underbrace{b_{T,x_1*x_2}}_{b_{T_1,x_1}\otimes b_{T_{\unsure{1|2}},x_2}}=
\]
\centerline{---}
% lines 18–25, formulary
\uline{Formulaire récapitulatif} pour ce formalisme
\[
(x_1,x_2)\longrightarrow x_1*x_2:\quad E_1\times E_2\longrightarrow E_1*E_2,
\]
où $E_i$ \inserted{$(i=1,2)$} est l'espace affine enveloppant d'un torseur $P_i$ sous un module \deleted{$\unsure{M_i}$} $L_i$, torseur muni d'un hom de $\chi$-trivialisation \inserted{($\chi\in\Gamma(X,O)$, donné)}
\[
\pi_i:P_i\longrightarrow L_i
\]
satisfaisant (par définition de ``$\chi$ trivialisation'')
\[
\pi_i(x_i+u_i)=\pi_i(x_i)+\chi\,u_i\qquad x_i\in\Gamma P_i,\ u_i\in\Gamma L_i
\]
% archive sheet number: 72
\scanpage{74}
% lines 1–5, linear extension
On peut aussi considérer l'extension linéaire \unsure{de|des} $\pi_i$ :
\[
E_{\unsure{j|i}}\longrightarrow L_i
\]
et la condition précédente signifie simplement que
\[
\pi_i(\deleted{\unsure{x}}u_i)=\unsure{\chi|\gamma}\cdot u_i\qquad\text{si }u_i\in L_i\subset E_i.
\]
% lines 6–10, fundamental formulas and right-hand glosses
Ceci posé, on aura les formules fondamentales qui suivent.

\textcircled{1} \uuline{bilinéarité} de $x_1*x_2$,
\margin{i.e. correspond à \uline{appl. lin.} $\boxed{E_1\otimes E_2\to E}$}
\[
\boxed{\psi(x_1*x_2)=\psi_1(x_1)\psi_{\unsure{1|2}}(x_2)}\tag*{\textcircled{2}}
\]
\margin{compatibilité aux augmentations}
où
\[
\begin{aligned}
E_1&\overset{\psi_1}{\mathrel{\text{---}}}\underline O\\
E_2&\overset{\psi_2}{\mathrel{\text{---}}}\underline O\\
E=E_1*E_2&\xrightarrow{\psi}\underline O
\end{aligned}
\]
% lines 11–15, augmentation fibres
sont les augmentations. En particulier, $*$ induit
\[
\begin{aligned}
P_1\times P_2&\longrightarrow P&\quad&\text{où}&P_i&=\psi_i^{-1}(\{1\})\\
x_1,x_2&\longmapsto x_1*x_2&&&P&=\psi^{-1}(\{1\})\\[-2pt]
&\text{induit par }*
\end{aligned}
\]
% lines 16–20, translation formulas, left annotation
\[
\boxed{\begin{aligned}
(x_1+u_1)*x_2&=x_1*x_2+\deleted{\unsure{\pi_2}}u_1\otimes\pi_2(x_2)\\
x_1*(x_2+u_2)&=x_1*x_2+\pi_1(x_1)\otimes u_2
\end{aligned}}
\]
\[
x_i\in\Gamma E_{\unsure{j|i}},\qquad u_i\in\Gamma L_i
\]
% slanted
\margin{\textcircled{3} Ces formules \unsure{utiles} sont celles \unsure{pour} \unsure{surtout} en se \unsure{restreignant} à des arguments $x_1\in\Gamma P_1,\ x_2\in\Gamma P_2$}
d'où aussitôt
\[
\boxed{\begin{aligned}
(x_1+u_1)*(x_2+u_2)={}&x_1*x_2+u_1\otimes\pi_2(x_2)+\pi_1(x_{\unsure{1|0}})\otimes u_2\\
&+\chi\,u_1\otimes u_2
\end{aligned}}
\]
% lines 21–26, restrictions at foot
Ces formules résultent de la bilinéarité $+$
\[
\boxed{\begin{aligned}
u_1*x_2&=u_1\otimes\pi_2(x_2)\\
x_1*u_2&=\pi_1(x_1)\otimes u_2
\end{aligned}}
\]
\deleted{ce qui} dont chacune implique (en faisant $x_2=u_2$ resp $x_1=u_1$)
\[
\boxed{u_1*u_2=\chi\,u_1\otimes u_2},
\]
compte tenu de $\pi_1(u_1)=\chi u_1,\quad\pi_2(u_2)=\chi(u_2)$
% archive sheet number: 73
\scanpage{75}
% lines 1–2, heading 4 and upper-left note
\textcircled{4} Relation de $\pi_i,T_i$
\[
\boxed{\pi_i(x)=\unsure{\chi|\gamma}x-T_i}
\]
% slanted
\margin{\unsure{remonte|remarque} aux généralités sur les $\chi$-trivialisations, n'a rien à voir avec l'op. $*$}
% lines 3–9, heading 5 and check on a submodule
\textcircled{5} Description de la $\chi$-trivialisation de $E$
\[
\boxed{\pi_E(x_1*x_2)=\pi_1(x_1)\otimes\pi_2(x_2)}
\]

Si $x_2=u_2\in\Gamma L_2$, on a $\inserted{x\overset{\mathrm{def}}{=}x_1*x_2={}}x_1*u_2=\pi_1(x_1)\otimes u_2\deleted{\unsure{\in E}}\inserted{\in\Gamma L_1\otimes L_2}$ et \deleted{on aura} le premier mb doit donc être égal à $\chi x=\chi\,\pi_1(x_1)\otimes u_2$, le deuxième est $\pi_1(x_1)\otimes\chi u_2$ puisque $\pi_2(u_2)=\chi u_2$, et c'est OK. \unsure{Item|Idem} pour $x_1=u_1$, $x_2$ quelcon.
% lines 10–15, corollary
Cor. On aura, pour $x_1\in\Gamma P_1$, $x_2\in\Gamma P_2$
\[
\underset{\substack{\shortparallel\\\chi x_1\otimes x_2-T}}{\pi_E(x_1*x_2)}
=\underbrace{(\chi x_1-T_1)}_{b_1}\otimes\underbrace{(\chi x_2-T_2)}_{b_2}
\]
donc
\[
\boxed{T=\chi(\,x_1\mathbin{\unsure{*|\otimes}}x_2)-b_1\otimes b_2}
\]
où
\[
b_i=\pi_{E^i}(x_i)=\chi x_i-T_i
\]
% lines 16–21, long slanted note at right
% slanted
\margin{NB La donnée de $x_1,x_2$ d'où $x_1*x_2$ définit \unsure{des} scindages\\
$E_1\simeq L_1\oplus\underline O\quad(u_1,\lambda_1)\mapsto\lambda_1 x_1+u_1$\\
$E_2\simeq L_2\oplus\underline O\quad(u_2,\lambda_2)\to\lambda_2 x_2+u_2$\\
$E\simeq\underbrace{L_1\otimes L_2}_{L}\oplus\unsure{U|O},\quad(u,\lambda)\mathbin{-}\lambda x_1*x_2+\deleted{\unsure{x}}u$\\
On aura alors\\
$(u_1,\lambda_1)\deleted{\unsure{\otimes}}*(u_2,\lambda_2)=$\\
$(\lambda_1\deleted{\unsure{b_2}}b_1\otimes u_2+u_1\otimes\lambda_2 b_2+\chi u_1\otimes u_2,\lambda_1\lambda_2)$}
% lines 22–26, verification at foot
Vérifions
\[
\boxed{\chi T=\underset{\substack{\shortparallel\\\chi x_1-b_1}}{T_1}*
\underset{\substack{\shortparallel\\\chi x_2-b_2}}{T_2}}
\]
le $2^{\mathrm{\unsure{d|o|e}}}$ \unsure{mb} est
\[
\begin{aligned}
&\chi^2 x_1*x_2-\chi\,\underbrace{\pi_1(x_1)}_{(\chi x_1-T_1)=b_1}\otimes b_2
-\chi\,b_1\otimes\underbrace{\pi_2(x_2)}_{(\chi x_2-T_2)=b_2}
+\underbrace{b_1*b_2}_{\chi b_1\otimes b_2}\\
&=\chi^2 x_1*x_2-\chi\,b_1\otimes b_2=\chi(\overbrace{\chi x_1*x_2-b_1\otimes b_2}^{T})
\end{aligned}
\]
c'est donc bien \deleted{ce qui avait \unsure{précédé|par $\chi$ de}} $\chi T$, OK
% archive sheet number: 74
\scanpage{76}
% lines 1–2, heading and superscribed qualification
$P=P_1*P_2$ \inserted{($P_i$ torseur $\unsure{\chi|\gamma}$-trivialisé, sous module $L_i$)} comme solution d'un pb universel
% lines 3–8, objects of the problem
On s'intéresse aux torseurs $\Pi$ sous un module $M$, et aux applications
\[
P_1\times P_2\xrightarrow{f}\Pi
\]
(qu'on va appeler ``biaffines'' compatibles avec les $\chi$-trivialisations données $\pi_1,\pi_2$ de $P_1,P_2$) qui satisfont la condition suivante :
% lines 9–14, bilinear datum and translations
\[
\exists\quad L_1\times L_2\xrightarrow{\alpha}\deleted{\unsure{\Pi}}M,\quad\text{application bilinéaire,}\tag*{$(*)$}
\]
telle qu'on ait
\[
\left\{\begin{aligned}
\deleted{\alpha}f(x_1+u_1,x_2)&=\alpha(x_1,x_2)+\deleted{B}\alpha(u_1,\pi_2(x_2))\\
\deleted{\alpha}f(x_1,x_2+u_2)&=\alpha(x_1,x_2)+\deleted{B}\alpha(\pi_1(x_1),u_2)
\end{aligned}\right.
\]
(d'où
\[
\begin{aligned}
\deleted{\alpha}f(x_1+u_1,x_2+u_2)={}&\alpha(x_1,x_2)+\deleted{B}\alpha(u_1,\pi_2(x_2))+
\deleted{B}\alpha(\pi_1(x_1),\unsure{x_2|u_2})\\
&+\chi\,\deleted{B}\alpha(u_1,u_2)
\end{aligned}
\]
% lines 15–21, dependence and qualification about uniqueness
Cela signifie donc que $\alpha$ est connu \inserted{(\unsure{en termes de|au lieu de|uniquement de} $B$)} quand on connait sa valeur sur \uline{un} couple $(x_1,x_2)$ de sections, et qu'il est donné par la formule précédente.

\deleted{On trouve que $B$ est} En fait, \deleted{$B$}$\alpha$ n'est pas néc. unique :
(car $(E_1\otimes L_2)\oplus(L_1\otimes E_2)\longrightarrow L_1\otimes L_2$ n'est pas néc. surjectif), et on va définir un \deleted{\unsure{Bi}} ``morphisme biaffine'' comme un couple $(f,\alpha)$, satisfaisant la condition de variance ci-dessus.
% archive sheet number: 75
\scanpage{77}
% lines 1–4, universal pair
Ceci posé,
\[
\bigl(\deleted{E}P_1*P_2,\quad P_1\mathbin{\unsure{*|\times}}P_2\xrightarrow{*}P_1*P_2\bigr)
\]
\deleted{\unsure{avec}} est une solution du pb universel défini par cette notion (pour $M,\Pi$ variables...).
% lines 5–14, compatible trivializations
Si on veut de plus décrire $P_1*P_2$ comme un torseur $\unsure{\chi|\gamma}$-trivialisé, on \deleted{définira} exigera, dans la notion de ``morphisme biaffine compatible avec les $\chi$-trivialisations'' (non seulement de $P_1,P_2$, mais aussi de $\Pi$), que l'on ait
\[
\pi_{\Pi}(\deleted{\unsure{\alpha}}f(x_1,x_2))=\pi_1(x_1)\otimes\pi_2(x_2)
\]
et on trouve maintenant une solution \unsure{de|du} ce problème universel pour le torseur $\chi$-trivialisé $P_1*P_2$.
% lines 15–20, morphisms
Mais on avait oublié d'expliciter la notion de morphisme de torseurs $\chi$-trivialisés,
\[
\begin{aligned}
P&\xrightarrow{f}P'\\
M&\xrightarrow{f_0}M'
\end{aligned}
\]
On exige
\[
\left\{\begin{aligned}
f(\overset{\substack{\Gamma P\\\rotatebox{-90}{$\in$}}}{x}+\overset{\substack{\Gamma M\\\rotatebox{-90}{$\in$}}}{u})&=f(x)+f_0(u)\\
\pi_{P'}(f(x))&=f_0\bigl(\pi_P(x)\bigr)
\end{aligned}\right.
\]
% archive sheet number: 76
\scanpage{78}
% lines 1–5, extensions and morphisms
Quand on interprète $P,P'$ en termes d'ext. de $\underline O$ par des Modules $L,L'$, munies de $T,T'$ \unsure{telle|telles} que $\psi(T)=\unsure{\chi|\gamma}$, $\psi'(T')=\chi$, la condition de compatibilité avec les $\chi$-trivialisations s'écrit simplement
\[
f(T)=T'.
\]
% lines 6–11, chosen splitting
\uline{Cas scindé} Un scindage d'un torseur $\chi$-trivialisé $(P,\deleted{\unsure{T}}\pi)$ est, par déf., une section (i.e. trivialisation) \inserted{$x_0$} de $P$. Alors \deleted{\unsure{à}}
\[
P\xleftarrow{\sim} L,\qquad u\longmapsto x_0+u
\]
\deleted{et la donnée de la $\pi$-tri} d'où
\[
E\xleftarrow{\sim}L+\underline O,\qquad(u,\lambda)\longmapsto\lambda x_0+u
\]
% lines 12–18, base point and coordinates
La donnée de la $\chi$-trivialisation équivaut à celle de l'élément
\[
\underset{\text{ou }b_{x_0}}{b}=\pi(x_0)=\chi x_0-T
\]
de \deleted{\unsure{terme}} sorte qu'en termes de $b$ on aura
\[
T=\chi x_0-b
\]
ou en termes de l'iso $E\simeq L+\underline O$,
\[
T=(-b,\chi)
\]
% lines 19–23, linear formula and check
On aura donc, en termes de $b$
\[
\pi(\overbrace{\underbrace{\lambda x_0+u}_{(u,\lambda)}}^{x})=\lambda b+\chi u
\]
[Vérifions que c'est \unsure{aussi|vrai|OK} \deleted{$\unsure{\pi}$} $\chi x-\psi(x)T=\chi\lambda x_0+\chi u-\lambda\underbrace{T}_{\chi x_0-b}$ \deleted{\unsure{$=$}}, on trouve bien $\chi u+\lambda b$ OK ]
% archive sheet number: 77
\scanpage{79}
% lines 1–6, split product
Si $P_i$ sous $L_i$, avec $\pi_i$, est scindé par $x_i$ $(i=1,2)$, alors $P_1*P_2$ est scindé par $x_1*x_2$. Sa $\unsure{\chi|\gamma}$-trivialisation est donnée par un élément $b$ de $\deleted{\unsure{M}}L_1\otimes L_2$, et on aura
\[
b=b_1\otimes b_2
\]
Il n'y a pas plus simple !
% lines 7–9, introduction to structural data
Ceci permet d'expliciter les données d'associativité, de commutativité, d'unité pour l'opération $*$. Mais de façon intrinsèque :
% lines 10–14, commutativity
1) \uline{Commutativité} L'isom \deleted{\unsure{d}}can. de commutativité
\[
\deleted{\unsure{P}}E_1*E_2\xrightarrow{\sim}E_2*E_1
\]
transforme $x_1*x_2$ en $x_2*x_1$, et $u_1\otimes u_2$ en $u_2\otimes u_1$, ce qui suffit à le caractériser.
% lines 15–20, associativity and beginning of unit
2) \uline{Associativité} L'iso d'ass.
\[
(E_1*E_2)*\deleted{\unsure{(}}E_3\xrightarrow{\sim}E_1*(E_2*E_3)
\]
transforme
\[
(x_1*x_2)*x_3\quad\text{en}\quad x_1*(x_2*x_3)
\]
et
\[
(\deleted{\unsure{x}}u_1\otimes u_2)\otimes u_3\quad\text{en}\quad u_1\otimes(u_{\unsure{2|1}}\otimes u_3),
\]
ce qui suffit à le caractériser.

3) \uline{Unité} : On \deleted{appelle} dit qu'une
% archive sheet number: 78
\scanpage{80}
% lines 1–5, qualification at upper edge and unit condition
\inserted{scindage $x_0$ d'un $P_0$ \unsure{muni déjà|muni} d'une} $\unsure{\chi|\gamma}$-trivialisation est ``\uline{unitaire}'', si l'élém
\[
b_0=\pi(x_0)
\]
de $L$ est \deleted{\unsure{déf}} \unsure{tel|telle} que $\{b\}$ soit une \uline{base} de $L$.
\deleted{Donc} Donc grâce à cette base, $L\xleftarrow{\sim}\underline O_X,\ b=1$, et on aura donc
% lines 6–10, coordinates for the unit
\[
\begin{aligned}
P_0&\simeq\underline O_{\deleted{X}},\quad x_0=1\\
E_0&=\underline O\times\underline O=\{(u,\lambda)\mid u,\lambda\in\Gamma O\}\xrightarrow{\psi}\lambda
\end{aligned}
\]
et
\[
\left|\begin{aligned}
\pi(u,\lambda)&=\lambda+\chi u\\
T&=(-1,\chi)
\end{aligned}\right.
\]
% lines 11–15, unit isomorphisms
Ceci posé, si $\deleted{E}P_1$ est un torseur $\chi$-trivialisé, alors les
\[
\left.\begin{aligned}
P_0*P_1&\simeq P_1\\
\underset{\shortparallel\,\underline O}{L_0}\otimes L_1&\simeq L_1\\
x_1&\longmapsto x_0*x_1\\
u_1&\mathbin{\unsure{\mapsto|-}}b_0\otimes\deleted{\unsure{x}}u_1
\end{aligned}\right.
\]
% lines 16–18, compatibility sentence
Les compatibilités entre les contraintes d'unité, d'assoc., de commutativité... .
% lines margin beside lines 11–20, diagonal writing from lower left
% slanted
\margin{\uline{NB} Un $(E,\psi,T)$ est \uline{inversible} pour $*$ sss\\
a) $L$ inversible\\
b) Sur \inserted{$X_0={}$}$V(\chi)\subset X$ (le schéma des zéros de $\chi$), la section $T_0=T|X_0$ de $({}_{\chi}P)_0\simeq{}_{\chi_0}P_0\simeq L_0$ est inversible. \deleted{\unsure{de}}\\
Si $(E,\psi,T)$ est scindé, et si $b\in\Gamma L$ correspond à un scindage $x_0$, \unsure{cf. devant|cf. dessus|alors} (on a $T=2\chi x-b$) donc $T_0=-b_0$, et la condition est que $b_0$ soit section inv de $L_0$ i.e. \unsure{qu’il en} soit une base. Cela équivaut à la condition $\unsure{L^{\otimes2}=2L^{\otimes2}+b\cdot L}$ \unsure{, ou à} $L',\unsure{L^{-\otimes}|L^{-1}}$ \deleted{\unsure{si}}\\
La condition d'inversibilité équivaut aussi à celle-ci\\
a) $L$ inversible\\
b) $\exists$ \uline{localement} un scindage $x$ de $(E,\psi,T)$, tel que le $b$ correspondant soit une base de $L$.}
% lines 19–21, new heading and specialization
\centerline{---}
Relation avec les involutions canoniques.

On suppose à partir de maintenant
\[
\chi=2.
\]
% archive sheet number: 79
\scanpage{81}
% lines 1–3
Soit $E$
\[
0\mathrel{\text{---}}L\mathrel{\text{---}}E\overset{\psi}{\mathrel{\text{---}}}O\mathrel{\text{---}}0
\]
une extension $2$-trivialisée par $T\in\Gamma E$ telle que $\psi(T)=2$, on va définir un automorphisme de la structure, par
% lines 4–8
\[
\sigma_E:\quad x\longmapsto\bar{x}=T\psi(x)-x
\qquad\left(=T-x\ \text{si }
  \begin{aligned}x&\in\Gamma P_1\big\}\\[-2pt]
  &\text{i.e. }\psi(x)=1\end{aligned}\right)
\]
On aura
\[
\psi(\bar{x})=\underbrace{\deleted{\unsure{2}}\psi(T)}_{2}\psi(x)-\psi(x)=\psi(x)
\]
i.e. compatibilité avec augmentation $\psi$.
% lines 9–14
On a bien sûr
\[
\overline{\bar{x}}=x\qquad\text{i.e.}\quad\sigma^2=\mathrm{id}_E,
\]
d'autre part
\[
\bar{\deleted{\unsure{x}}u}=-u\qquad\text{si }u\in L
\]
donc $\sigma$ induit $-\mathrm{id}_L$ sur $L$, $\mathrm{id}_O$ sur $\underline O\simeq E/L$.
% lines 15–18
\uline{Si $2$ inversible}\inserted{dans $O$, plus gén. si $2\,\mathrm{id}_M$ est inv.,}, alors on a un scindage canonique
\[
\underline E\simeq L+\underset{\substack{\cap\\E}}{\deleted{M}E_0},
\]
\margin{NB. L'élément $x_0\in\Gamma E_0$ tel que $\psi(x_0)=1$ est caractérisé par $2x_0=T$ i.e. $x_0=\frac{1}{2}T$}
% lines 19–23
où
\[
\underset{\substack{\sim\\[-2pt]\psi|E_0\ \underline O_X}}{E_0}=E^\sigma=
\underset{\substack{\shortparallel\\\operatorname{Im}(1+\sigma)\\\shortparallel\\
\{x+\bar{x}\mid x\in E\}}}{\operatorname{Ker}(1-\sigma)}
=\{x\in E\mid x=\bar{x}\}
\]
\deleted{Donc on \unsure{peut}}

\deleted{\unsure{Soit}}On aura donc
\[
(u,\lambda)^{\rule{.62em}{.35pt}}=(-u,\lambda)
\]
% lines 24–26
et
\[
T=\deleted{\unsure{(0,2)\psi(x)}}x+\bar{x}\qquad\text{pour }\psi(x)=1
\]
d'où
\[
\boxed{T=(0,2)}\qquad\text{On aura}\quad\boxed{\pi(u,\lambda)=(2u,0)}
\]
% archive sheet number: 80
